\begin{lemma}
If $X$ is infinite, then $\mathsf{CardinalityFront}(k,X)=[X]^k$ is a
front on $X$ for every $k\in\mathbb N$.
\end{lemma}
\begin{proof}
Put $F=\mathsf{CardinalityFront}(k,X)$, as defined in
\noderef{CardinalityFront}.  We verify the four fields in the explicit
definition of a front in \noderef{Front}.  The infinite-base field is
exactly the given hypothesis $h_X:X$ is infinite.

We first verify the base alternative.  Suppose that $k=0$.  By the
definition of $F$, a finite set $s$ belongs to $F$ exactly when

\[
|s|=0\quad\text{and}\quad s\subseteq X.
\]

The first condition is equivalent to $s=\emptyset$.  The empty set does
satisfy the second condition, so in this case

\[
F=\{\emptyset\}.
\tag{1}
\]

Now suppose that $k>0$.  We prove
$\mathsf{FrontBase}(F)=X$, using the definition of the base in
\noderef{FrontBase}.  If $n\in\mathsf{FrontBase}(F)$, there is an
$s\in F$ with $n\in s$.  Membership $s\in F$ gives $s\subseteq X$ by
\noderef{CardinalityFront}, and hence $n\in X$.  This proves
$\mathsf{FrontBase}(F)\subseteq X$.

For the converse, fix $x\in X$.  The set $X\setminus\{x\}$ is infinite:
if it were finite, then adjoining the at most one point $x$ would make
$X$ finite, contrary to $h_X$.  An infinite set contains any prescribed
finite number of pairwise distinct points.  Indeed, choose them
inductively; after finitely many choices, their finite set cannot exhaust
an infinite set, so another point remains.  Therefore choose pairwise
distinct

\[
a_1,\ldots,a_{k-1}\in X\setminus\{x\};
\]

when $k=1$ this is the empty list.  Let
$s=\{x,a_1,\ldots,a_{k-1}\}$.  All its displayed elements are distinct,
so $|s|=k$, and every element of $s$ lies in $X$.  Thus
$s\in\mathsf{CardinalityFront}(k,X)$, and $x\in s$.  The definition of
\noderef{FrontBase} now gives
$x\in\mathsf{FrontBase}(F)$.  Hence $X\subseteq\mathsf{FrontBase}(F)$,
which proves

\[
\mathsf{FrontBase}(F)=X\qquad(k>0).
\tag{2}
\]

Equations (1) and (2) establish the base alternative in all cases.

Next let $s,t\in F$ and assume
$\mathsf{InitialSegment}(s,t)$.  If $s=t$, the required conclusion is
immediate.  Otherwise, the non-equality clause in the definition of
\noderef{InitialSegment} supplies an $n\in t$ such that

\[
s=\{m\in t\mid m<n\}.
\]

Thus $s\subseteq t$, while $n\notin s$ because $n<n$ is false.  Hence
$s$ is a proper subset of the finite set $t$, so $|s|<|t|$.  On the
other hand, membership of both sets in $F$, by
\noderef{CardinalityFront}, gives $|s|=k=|t|$, a contradiction.  Thus
the prefix-free clause from \noderef{Front} holds.

It remains to prove density.  Let $Y\subseteq X$ be infinite.  By
\noderef{InfiniteSetEnumeration}, choose an increasing enumeration
$y\in\mathsf{IncSeq}$ such that
$\operatorname{range}(y)=Y$.  By the definition of \noderef{IncSeq},
$y$ is an order embedding: it is injective, it satisfies
$i<j\Rightarrow y(i)<y(j)$, and it reflects strict order.  The range
equality also says that every $y(i)$ belongs to $Y$ and that every member
of $Y$ is $y(j)$ for some $j$.

First suppose that $k=0$.  Take $s=\emptyset$.  Then $s\in F$ by the
definition of $F$.  The element $y(0)$ belongs to $Y$.  If $m\in Y$ and
$m<y(0)$, write $m=y(j)$ using the range equality.  Order reflection
would give $j<0$, which is impossible.  Consequently

\[
\forall m,\qquad m\in\emptyset
\quad\Longleftrightarrow\quad m\in Y\text{ and }m<y(0).
\]

Together with $y(0)\in Y$, this is exactly
$\mathsf{ProperPrefixSet}(\emptyset,Y)$ by the definition in
\noderef{ProperPrefixSet}.  This supplies the required density witness in
the case $k=0$.

Now suppose that $k>0$ and let

\[
s=\{y(i)\mid i<k\}.
\]

The values $y(0),\ldots,y(k-1)$ are pairwise distinct by injectivity, so
$|s|=k$.  Each value belongs to $Y\subseteq X$, hence $s\subseteq X$.
Therefore $s\in F$.  We use $y(k)$ as the cutoff.  It belongs to $Y$,
and for every natural number $m$ we have

\[
m\in s
\quad\Longleftrightarrow\quad m\in Y\text{ and }m<y(k).
\tag{3}
\]

For the forward implication, write $m=y(i)$ with $i<k$; strict increase
gives $y(i)<y(k)$.  For the reverse implication, write $m=y(j)$ using
$m\in Y$.  From $y(j)=m<y(k)$ and order reflection we get $j<k$, so
$m=y(j)\in s$.  Equation (3), together with $y(k)\in Y$, is precisely
$\mathsf{ProperPrefixSet}(s,Y)$ by \noderef{ProperPrefixSet}.

We have therefore established the infinite-base field, the base
alternative, prefix-freeness, and density.  By the four defining fields
of \noderef{Front}, these facts prove
$\mathsf{Front}(\mathsf{CardinalityFront}(k,X),X)$.
\end{proof}
